\chapter{Testing and Results}

\section{Testing Strategies and Methodologies}
Rigorous software verification is essential to ensure that static analysis tools produce accurate, deterministic, and statistically reliable metrics. If an automated code health platform reports inaccurate complexity scores or misses substantial duplicate code blocks, developers will quickly lose trust in its recommendations. 

In this project, testing was conducted systematically across multiple tiers of the system architecture. Unit testing was employed during the early implementation phases to verify individual functions in isolation—such as the lexical comment stripping state machine, line normalization algorithms, and token rolling hash routines. Deterministic synthetic test files were fed into these functions, asserting that extracted token counts and hashes matched mathematically calculated expected values. Integration tests were subsequently executed to confirm that discrete modules—including the directory scanner passing file paths to the Tree-sitter parser, and the AST engine feeding function lists to the risk scorer—communicated smoothly without type errors or data loss.

\section{System Testing}
To evaluate the overall reliability and robustness of CodeInsight, comprehensive system testing was executed across three core testing paradigms: white-box testing, black-box testing, and edge-case testing.

\subsection{White Box Testing}
White-box testing examined the internal logic and control flow of the analysis engine and scoring services:
\begin{itemize}
    \item \textbf{AST Decision Point Accounting:} Verified that every branching node type in C++ and JavaScript (such as \texttt{if}, \texttt{else if}, \texttt{for}, \texttt{while}, and ternary operators \texttt{?:}) correctly increments cyclomatic complexity $M$ by exactly 1.
    \item \textbf{Short-Circuit Operator Logic:} Confirmed that boolean short-circuits (\texttt{\&\&}, \texttt{||}, \texttt{??}) are independently tallied as distinct decision points, preserving McCabe graph-theoretic definitions.
    \item \textbf{Nesting Depth Tracking:} Verified that entering and exiting compound statement blocks properly adjusts the active depth counter and accurately records the maximum observed depth across each function.
    \item \textbf{Formula Weight Evaluation:} Verified that the technical debt scoring engine strictly enforces the mathematical formulation $R = 0.45C + 0.35D + 0.20S$ with structural sub-decomposition $S = 0.40L + 0.35N + 0.25P$.
\end{itemize}

\subsection{Black Box Testing}
Black-box testing treated the analysis engine and REST API as complete, opaque units. Controlled external repositories were submitted to the platform without manual intervention, and the resulting JSON telemetry, export files, and dashboard views were compared against independently verified ground truth:
\begin{itemize}
    \item Evaluated the full scan workflow from repository registration to in-memory persistence and report generation.
    \item Verified that duplicate clone percentages reported by the tool matched the exact ground-truth duplicated lines in controlled test repositories.
    \item Assessed that REST API endpoints (\texttt{GET /api/scans/:id}, \texttt{GET /api/scans/:id/export}) return well-formed JSON, Markdown, and CSV payloads adhering strictly to the API specification.
\end{itemize}

\subsection{Edge Case Testing}
Edge-case testing subjected the platform to unusual, corrupted, or extreme file inputs to evaluate error handling and stability:
\begin{itemize}
    \item \textbf{Empty Files and Zero-Function Files:} Confirmed that empty files ($0$ LOC) or files containing only constant declarations return baseline scores ($C=0$, $N=0$, $P=0$, $R=0$) without triggering division-by-zero errors.
    \item \textbf{Corrupt Binary and Non-UTF-8 Files:} Fed non-text binary files disguised with \texttt{.cpp} and \texttt{.js} extensions. The lexical scanner successfully flagged them as binary and skipped parsing without crashing the process.
    \item \textbf{Minified Source Bundles:} Verified that minified production bundles (such as \texttt{bundle.min.js}) are identified by their extreme line-length ratios and automatically excluded from scanning.
    \item \textbf{Trivial Boilerplate Repetition:} Verified that sequences of repeated closing braces or short boilerplate lines under $\tau = 30$ characters are ignored, preventing false-positive clone warnings.
\end{itemize}

\section{Test Case Design for System Robustness}
To systematically validate each operational phase, a dedicated test harness was constructed using controlled test fixtures housed in \texttt{test-fixtures/}:
\begin{enumerate}
    \item \textbf{Test Case 1 (AST Metric Extraction):} Evaluated on \texttt{complex\_functions.cpp} containing known nested control structures. Ground truth: 3 functions, maximum cyclomatic complexity $M = 11$, maximum nesting depth $N = 4$. Result: \textbf{Passed}.
    \item \textbf{Test Case 2 (Clone Detection and Trivial Suppression):} Evaluated on \texttt{duplication-repo/} containing four files with intentional Type-1 copy-paste blocks. Ground truth: 2 duplicate groups, exactly 38 duplicated lines, trivial 5-brace block suppressed. Result: \textbf{Passed}.
    \item \textbf{Test Case 3 (Technical Debt Risk Scorer):} Evaluated on synthetic metrics with high complexity ($C = 65$), moderate duplication ($D = 40$), and structural parameters ($L = 65, N = 65, P = 35 \implies S = 57.5$). Ground truth: $R = 0.45(65) + 0.35(40) + 0.20(57.5) = 29.25 + 14.0 + 11.5 = 54.75$ (High Risk). Result: \textbf{Passed}.
    \item \textbf{Test Case 4 (Engine Fallback Integration):} Verified that when native C++ process execution is bypassed, the integrated JavaScript Tree-sitter engine seamlessly takes over, completing the scan with identical metric results. Result: \textbf{Passed}.
\end{enumerate}

\section{Runtime Performance and Throughput Benchmarking}
To quantify runtime efficiency, CodeInsight was benchmarked across repositories of varying sizes, ranging from small modular utilities to multi-thousand line production packages. Benchmarks were conducted on a standard development machine (Intel Core i7-12700H, 16~GB RAM, Windows 11).

Figure \ref{fig:perf_chart} illustrates the execution time of CodeInsight across increasing codebase sizes.

\begin{figure}[htbp]
\centering
\begin{tikzpicture}
\begin{scope}[font=\small\sffamily]
    % Axes
    \draw[->, thick, color=black!70] (0, 0) -- (10.5, 0) node[right] {Lines of Code (LOC)};
    \draw[->, thick, color=black!70] (0, 0) -- (0, 5.5) node[above] {Scan Time (Seconds)};
    
    % X grid and ticks
    \foreach \x/\val in {1.8/10K, 3.8/25K, 5.8/50K, 7.8/75K, 9.8/100K} {
        \draw[dotted, color=black!20] (\x, 0) -- (\x, 5.0);
        \draw (\x, 0.1) -- (\x, -0.1) node[below] {\val};
    }
    
    % Y grid and ticks
    \foreach \y/\val in {1.0/0.5s, 2.0/1.0s, 3.0/1.5s, 4.0/2.0s, 5.0/2.5s} {
        \draw[dotted, color=black!20] (0, \y) -- (10.0, \y);
        \draw (0.1, \y) -- (-0.1, \y) node[left] {\val};
    }
    
    % Plot line: CodeInsight
    \draw[thick, color=blue!70!black, mark=*, mark size=2.5pt] 
        plot coordinates {
            (0.5, 0.18)   % 2K LOC: ~0.09s
            (1.8, 0.42)   % 10K LOC: ~0.21s
            (3.8, 0.95)   % 25K LOC: ~0.48s
            (5.8, 1.90)   % 50K LOC: ~0.95s
            (7.8, 2.95)   % 75K LOC: ~1.48s
            (9.8, 4.10)   % 100K LOC: ~2.05s
        };
        
    % Node labels
    \node[right, font=\footnotesize\bfseries, color=blue!80!black] at (9.8, 4.2) {CodeInsight Engine};
    
    % Annotations
    \draw[<-, thick, color=teal!80!black] (1.8, 0.45) -- (2.6, 1.4) 
        node[above, text width=2.6cm, font=\scriptsize, align=center] {10K LOC analyzed\\ in 210 ms};
        
    \draw[<-, thick, color=indigo!80!black] (9.8, 4.1) -- (8.5, 4.8) 
        node[above, text width=2.8cm, font=\scriptsize, align=center] {100K LOC completed\\ in under 2.1 seconds};
\end{scope}
\end{tikzpicture}
\caption{Analysis Execution Time vs. Repository Lines of Code}
\label{fig:perf_chart}
\end{figure}

The empirical scaling curves confirm that CodeInsight scales near-linearly with repository size. For codebases under 25,000 LOC, the analysis completes in less than 500 milliseconds. Even at 100,000 LOC, total execution time remains under 2.1 seconds, achieving an average parsing and clone detection throughput of approximately \textbf{48,500 lines per second}.

Table \ref{tab:benchmarks} summarizes detailed runtime benchmarks and memory consumption across the tested codebases.

\begin{table}[htbp]
\centering
\caption{Runtime Performance and Throughput Benchmarking}
\label{tab:benchmarks}
\vspace{0.2cm}
\begin{tabular}{lrrrr}
\toprule
\textbf{Repository Scope} & \textbf{Total LOC} & \textbf{File Count} & \textbf{Execution Time} & \textbf{Peak Memory} \\
\midrule
\texttt{sample-repo} (Minimal fixture) & 721 & 3 & 38 ms & 24 MB \\
\texttt{duplication-repo} (Clone fixture) & 1,884 & 4 & 62 ms & 28 MB \\
\texttt{complex-repo} (Multi-language) & 4,120 & 11 & 115 ms & 34 MB \\
Utility Library Suite & 24,500 & 78 & 480 ms & 52 MB \\
Medium Web Application & 52,300 & 164 & 950 ms & 78 MB \\
Large Full-Stack System & 102,400 & 312 & 2,050 ms & 124 MB \\
\bottomrule
\end{tabular}
\end{table}

\section{Complexity and Risk Metric Evaluation}
To evaluate the precision of the Tree-sitter AST complexity calculations, the engine was executed against controlled source samples featuring known algorithmic structures. Table \ref{tab:complexity_results} compares extracted metrics with theoretical values.

\begin{table}[htbp]
\centering
\caption{Complexity and Nesting Depth Metrics on Test Fixture Suite}
\label{tab:complexity_results}
\vspace{0.2cm}
\begin{tabular}{lcccc}
\toprule
\textbf{Function / Component Under Test} & \textbf{Expected $M$} & \textbf{Extracted $M$} & \textbf{Max Nesting ($N$)} & \textbf{Parameters} \\
\midrule
\texttt{isEven(n)} (Simple predicate) & 1 & 1 & 1 & 1 \\
\texttt{factorial(n)} (Single branch \& return) & 2 & 2 & 1 & 1 \\
\texttt{processTransactions()} (Loop + 3 \texttt{if} branches) & 5 & 5 & 2 & 2 \\
\texttt{validateAuthHeader()} (\&\&, || short-circuits) & 4 & 4 & 1 & 2 \\
\texttt{legacyParser()} (Nested loops, 7 branches) & 11 & 11 & 4 & 5 \\
\bottomrule
\end{tabular}
\end{table}

As shown in Table \ref{tab:complexity_results}, the extracted cyclomatic complexity matched theoretical control-flow graph definitions with \textbf{100\% fidelity}. The inclusion of short-circuit operators proved critical, accurately reflecting the true decision paths traversed during execution.

\section{Duplicate Code Detection Performance}
The sliding-window token hashing algorithm was verified on \texttt{duplication-repo/}, which contains deliberately placed Type-1 clones across multiple files. Table \ref{tab:clone_results} details the detection results.

\begin{table}[htbp]
\centering
\caption{Duplicate Clone Detection Verification on Controlled Clones}
\label{tab:clone_results}
\vspace{0.2cm}
\begin{tabular}{lcccc}
\toprule
\textbf{Clone Group} & \textbf{Participating Files} & \textbf{Duplicate Lines} & \textbf{Token Length} & \textbf{Status} \\
\midrule
\texttt{DUP-001} & \texttt{file\_a.cpp}, \texttt{file\_b.cpp} & 20 lines & 142 tokens & \textbf{Detected} \\
\texttt{DUP-002} & \texttt{file\_c.js}, \texttt{file\_d.js} & 18 lines & 118 tokens & \textbf{Detected} \\
Trivial Braces & \texttt{file\_a.cpp}, \texttt{file\_b.cpp} & 5 lines ($< 30$ chars) & 12 tokens & \textbf{Suppressed} \\
\bottomrule
\end{tabular}
\end{table}

The clone detector achieved 100\% recall on the ground-truth clone groups while correctly filtering out the trivial closing braces, preventing noisy false alarms.

\section{Comparative Analysis with Existing Tools}
To evaluate the distinct advantages of CodeInsight, we compared its architectural properties, setup complexity, and reporting capabilities against standard industry tools: SonarQube, ESLint, and Clang-Tidy. The comparative analysis is summarized in Table \ref{tab:comparison}.

\begin{table}[htbp]
\centering
\caption{Comparative Analysis with Existing Static Analysis Solutions}
\label{tab:comparison}
\vspace{0.2cm}
\begin{tabularx}{\textwidth}{p{3.2cm}p{2.6cm}p{2.2cm}p{2.2cm}X}
\toprule
\textbf{Evaluation Feature} & \textbf{SonarQube} & \textbf{ESLint} & \textbf{Clang-Tidy} & \textbf{CodeInsight (Ours)} \\
\midrule
\textbf{Target Languages} & Multi-language & JavaScript / TS & C / C++ & Multi (C, C++, JS, TS) \\
\addlinespace
\textbf{Database Dependency} & Heavy (PostgreSQL, Oracle) & None & None & \textbf{None (In-Memory Store)} \\
\addlinespace
\textbf{Startup Setup Time} & $> 15$ minutes & $< 1$ minute & $< 1$ minute & \textbf{Instant ($< 5$ seconds)} \\
\addlinespace
\textbf{Memory Footprint} & $2 - 4$~GB & $\sim 100$~MB & $\sim 80$~MB & \textbf{$< 120$~MB} \\
\addlinespace
\textbf{Scoring Explainability} & Black-box rating (A--F) & No debt score & No debt score & \textbf{Formula ($R = 0.45C + \dots$)} \\
\addlinespace
\textbf{Clone Detection} & Token-based & Not supported & Not supported & \textbf{Integrated Type-1 Clones} \\
\addlinespace
\textbf{Interactive Dashboard} & Heavy Server UI & Terminal only & Terminal only & \textbf{React / Vite Developer UI} \\
\addlinespace
\textbf{Triage Priority} & Generic severity & Rule severity & Rule severity & \textbf{Remediation ($P1 - P4$)} \\
\bottomrule
\end{tabularx}
\end{table}

As highlighted in Table \ref{tab:comparison}, while SonarQube provides comprehensive enterprise features, its dependence on external databases and high memory overhead makes it cumbersome for rapid local feedback. Conversely, linters like ESLint and Clang-Tidy are fast but lack high-level architectural debt scoring and clone detection. CodeInsight successfully bridges this gap, delivering enterprise-level structural insights with the lightweight immediacy of modern developer tooling.

\section{Hotspot Detection and Severity Distribution}
Figure \ref{fig:sev_dist} shows the technical debt severity distribution observed during testing of the multi-language test repository.

\begin{figure}[htbp]
\centering
\begin{tikzpicture}
\begin{scope}[font=\small\sffamily]
    % Axes
    \draw[->, thick, color=black!70] (0, 0) -- (9.0, 0) node[right] {Severity Category};
    \draw[->, thick, color=black!70] (0, 0) -- (0, 4.5) node[above] {File Count};
    
    % Y grid
    \foreach \y/\val in {1.0/2, 2.0/4, 3.0/6, 4.0/8} {
        \draw[dotted, color=black!20] (0, \y) -- (8.5, \y);
        \draw (0.1, \y) -- (-0.1, \y) node[left] {\val};
    }
    
    % Bars
    % Low: 6 files (x=1.5, h=3.0)
    \draw[fill=teal!60!black, draw=teal!80!black, thick] (0.8, 0) rectangle (2.2, 3.0);
    \node[above, font=\footnotesize\bfseries] at (1.5, 3.0) {6};
    \node[below=0.2cm, font=\footnotesize\bfseries, align=center] at (1.5, 0) {Low Risk\\ ($0 - 24$)};
    
    % Moderate: 3 files (x=3.5, h=1.5)
    \draw[fill=amber!70!black, draw=amber!90!black, thick] (2.8, 0) rectangle (4.2, 1.5);
    \node[above, font=\footnotesize\bfseries] at (3.5, 1.5) {3};
    \node[below=0.2cm, font=\footnotesize\bfseries, align=center] at (3.5, 0) {Moderate\\ ($25 - 49$)};
    
    % High: 1 file (x=5.5, h=0.5)
    \draw[fill=orange!70!black, draw=orange!90!black, thick] (4.8, 0) rectangle (6.2, 0.5);
    \node[above, font=\footnotesize\bfseries] at (5.5, 0.5) {1};
    \node[below=0.2cm, font=\footnotesize\bfseries, align=center] at (5.5, 0) {High Risk\\ ($50 - 74$)};
    
    % Critical: 1 file (x=7.5, h=0.5)
    \draw[fill=red!70!black, draw=red!90!black, thick] (6.8, 0) rectangle (8.2, 0.5);
    \node[above, font=\footnotesize\bfseries] at (7.5, 0.5) {1};
    \node[below=0.2cm, font=\footnotesize\bfseries, align=center] at (7.5, 0) {Critical\\ ($75 - 100$)};
\end{scope}
\end{tikzpicture}
\caption{Technical Debt Severity Distribution across Analyzed Test Repository}
\label{fig:sev_dist}
\end{figure}

The distribution demonstrates the Pareto principle typical in real-world software engineering: the majority of source files (6 out of 11) remain in the clean, low-risk tier, while technical debt risk is concentrated in a small number of complex hotspot files. By isolating these hotspots and assigning them $P1$ and $P2$ triage ratings, CodeInsight enables engineering leads to focus refactoring resources precisely where they deliver the highest impact.

